\subsection{Do neutrinos dream in 5D? Towards a comprehensive extra-dimensional neutrino phenomenology --- arXiv:2512.02101}

\textbf{Authors:} Arturo de Giorgi, Dhruv Pasari, Jessica Turner

\paragraph{主な主張}
$S^1/\mathbb{Z}_2$上の5次元ニュートリノ模型について、brane/bulkのDirac/Majorana質量を特徴とする4設定を比較し、質量項の種類と位置によってKKスペクトル・振動・実験制約が大きく変わることを示す\cite{deGiorgi:2026Neutrinos5D}。

\paragraph{新規性}
flavour質量行列を同時対角化できる仮定の下で4設定の質量・混合を解析的に整理し、MINOS/MINOS+・Daya Bayとの比較からbulk Majoranaの共鳴構造やbrane Majoranaで半径感度が弱まる領域を示した。

\paragraph{位置付け}
最小Dirac模型からMajoranaを含む拡張までを接続する体系的な5次元研究であり、$T^2$への拡張で固定している質量生成機構と境界条件を明確にするための基礎になる。
